% ============================================================
% REMERCIEMENTS
% ============================================================

\newpage
\thispagestyle{empty}

\begin{center}
{\Large\textbf{Remerciements}}
\end{center}

\vspace{1.5cm}

Au terme de ce projet de fin d'études, je tiens à exprimer ma profonde gratitude envers toutes les personnes qui ont contribué à la réalisation de ce travail.

\vspace{0.5cm}

Je remercie tout d'abord \textbf{\underline{\hspace{4cm}}}, mon encadrant académique à l'Institut Supérieur de Gestion de Tunis, pour ses conseils avisés, son suivi rigoureux et sa disponibilité tout au long de ce projet.

\vspace{0.5cm}

Je tiens également à remercier \textbf{\underline{\hspace{4cm}}}, mon encadrant professionnel au sein de \textbf{TIM Tunisie}, pour la confiance qu'il m'a accordée, son accompagnement technique et ses orientations précieuses dans la compréhension des besoins métier liés aux systèmes \textbf{QALITAS} et \textbf{GMAO PRO}.

\vspace{0.5cm}

Mes remerciements s'adressent aussi à l'ensemble de l'équipe de \textbf{TIM Tunisie} pour l'accueil chaleureux, l'environnement de travail stimulant et les échanges enrichissants qui ont facilité la réalisation de cette plateforme.

\vspace{0.5cm}

Je remercie sincèrement les membres du jury pour l'honneur qu'ils me font en acceptant d'évaluer ce travail. Leurs remarques et suggestions seront d'un apport considérable.

\vspace{0.5cm}

Enfin, je ne saurais oublier mes parents, ma famille et mes amis pour leur soutien moral indéfectible et leurs encouragements permanents.

\vspace{1cm}

\begin{flushright}
\textit{Houcine LAGHMANI}\\
\textit{Tunis, Juin 2026}
\end{flushright}

\newpage
